\chapter{Complexity — Universal K}
%%% generated by the blueprint pipeline from TCSlib.Complexity.TuringMachine.UniversalK
\section{Overview}

Declarations of the Lean module \texttt{TCSlib.Complexity.TuringMachine.UniversalK}.

\section{Declarations}

\begin{definition}[Effective coding for a fixed tape count]
\lean{Turing.EffectiveCodeK}
\label{Turing.EffectiveCodeK}
\leanok
\uses{Turing.FinTM, Turing.FinTM.ComputesFunInTime, Turing.MultiTapeTM, Turing.serializeK}
For a natural number $k$, an effective code consists of a function assigning each
Boolean string $\alpha$ a natural number $N(\alpha)$, a machine with $k$ work tapes and
state set $\mathrm{Fin}(N(\alpha)+1)$ for each $\alpha$, and a finite Boolean machine
called the canonizer. It also includes a natural-valued time bound for the canonizer and
a proof that the canonizer computes the serialization of the machine assigned to
$\alpha$ within that bound.
\end{definition}

\begin{definition}[Finite machine decoded from a codeword]
\lean{Turing.EffectiveCodeK.toFinTM}
\label{Turing.EffectiveCodeK.toFinTM}
\leanok
\uses{Turing.EffectiveCodeK, Turing.FinTM}
For a natural number $k$, an effective coding scheme for machines with $k$ work tapes,
and a finite Boolean word $\alpha$, the associated finite Turing machine has $k$ work
tapes, state set $\{0,1,\ldots,n\}$ where $n$ is the number assigned to $\alpha$ by the
coding scheme, and underlying machine equal to the machine decoded from $\alpha$.
\end{definition}

\begin{definition}[Effective coding from serialization]
\lean{Turing.EffectiveCodeK.ofSerialize}
\label{Turing.EffectiveCodeK.ofSerialize}
\leanok
\uses{Turing.EffectiveCodeK, Turing.FinTM, Turing.FinTM.ComputesFunInTime, Turing.MultiTapeTM, Turing.serializeK}
Let $k$ be a natural number, let $\mathrm{Code}$ be a type, and assign to each
$M\in\mathrm{Code}$ a natural number of states minus one and a $k$-work-tape machine
over Boolean symbols with that state set. Suppose a serialization of each code equals
the fixed serialization of its assigned machine, and a decoder maps every Boolean word
to a code. If a finite Boolean-alphabet Turing machine computes the serialization of the
decoded code in time bounded by $T$, then these data determine an effective coding
scheme for $k$-work-tape machines: each word decodes to its assigned machine, and the
same machine computes that machine’s fixed serialization within time $T$.
\end{definition}

\begin{definition}[One-tape effective coding]
\lean{Turing.EffectiveMachineCode.toK}
\label{Turing.EffectiveMachineCode.toK}
\leanok
\uses{Turing.CodeTM.serialize, Turing.EffectiveCodeK, Turing.EffectiveCodeK.ofSerialize, Turing.EffectiveMachineCode, Turing.oneTapeMachineFormat}
Every effective machine coding scheme whose decoded machines have one work tape gives an
effective coding scheme for one-work-tape machines, with each Boolean word decoding to
the same machine and with its canonizer computing the fixed serialization of that
machine.
\end{definition}

\begin{definition}[Two-tape effective coding]
\lean{Turing.EffectiveMachineCode2.toK}
\label{Turing.EffectiveMachineCode2.toK}
\leanok
\uses{Turing.Code2TM.serialize, Turing.Code2TM.serialize_eq_serializeK, Turing.EffectiveCodeK, Turing.EffectiveCodeK.ofSerialize, Turing.EffectiveMachineCode2}
Every effective machine coding scheme whose decoded machines have two work tapes gives
an effective coding scheme for two-work-tape machines, with each Boolean word decoding
to the same machine and with its canonizer computing the fixed serialization of that
machine.
\end{definition}

\begin{theorem}[Conversion to a Finite Machine]
\lean{Turing.EffectiveMachineCode.toK_toFinTM}
\label{Turing.EffectiveMachineCode.toK_toFinTM}
\leanok
\uses{Turing.CodeTM.toFinTM, Turing.EffectiveCodeK.toFinTM, Turing.EffectiveMachineCode, Turing.EffectiveMachineCode.toK}
\difficulty{1}
Let $c$ be an effective machine code and let $\alpha$ be a binary string. The finite
machine obtained by converting the $K$-machine associated with $c$ at $\alpha$ is equal
to the finite machine obtained by converting the machine decoded from $\alpha$.
\end{theorem}

\begin{theorem}[Conversion to a Finite Machine]
\lean{Turing.EffectiveMachineCode2.toK_toFinTM}
\label{Turing.EffectiveMachineCode2.toK_toFinTM}
\leanok
\uses{Turing.Code2TM.toFinTM, Turing.EffectiveCodeK.toFinTM, Turing.EffectiveMachineCode2, Turing.EffectiveMachineCode2.toK}
\difficulty{1}
Let $c$ be an effective two-work-tape machine code and let $\alpha$ be a binary string.
The finite machine obtained by converting the $K$-machine associated with $c$ at
$\alpha$ is equal to the finite machine obtained by converting the code decoded from
$\alpha$.
\end{theorem}

\begin{theorem}[Decoded Machine from a Binary String]
\lean{Turing.EffectiveMachineCode.toK_machine}
\label{Turing.EffectiveMachineCode.toK_machine}
\leanok
\uses{Turing.EffectiveMachineCode, Turing.EffectiveMachineCode.toK}
\difficulty{1}
Let $c$ be an effective machine code and let $\alpha$ be a binary string. The machine
assigned to $\alpha$ by the $K$-machine construction associated with $c$ is equal to the
underlying machine decoded from $\alpha$.
\end{theorem}

\begin{theorem}[Decoded Two-Tape Machine]
\lean{Turing.EffectiveMachineCode2.toK_machine}
\label{Turing.EffectiveMachineCode2.toK_machine}
\leanok
\uses{Turing.EffectiveMachineCode2, Turing.EffectiveMachineCode2.toK}
\difficulty{1}
Let $c$ be an effective two-work-tape machine code and let $\alpha$ be a binary string.
The two-work-tape machine assigned to $\alpha$ by the $K$-machine construction
associated with $c$ is equal to the underlying machine decoded from $\alpha$.
\end{theorem}

\begin{theorem}[Variable-Tape Self-Description Startup]
\lean{Turing.exists_selfStartupTMK}
\label{Turing.exists_selfStartupTMK}
\leanok
\uses{Turing.EffectiveCodeK, Turing.FinTM.bufferTape, Turing.MultiTapeTM, Turing.MultiTapeTM.initCfg, Turing.MultiTapeTM.runFrom, Turing.MultiTapeTM.spaceUsed, Turing.pairEncode, Turing.tableBitsK, Turing.unaryFin}
For every natural number $k$ and every effective code system $E$ assigning each binary
string $\alpha$ a deterministic binary machine with $k$ work tapes and a finite indexed
state set, there are a natural number $K$, a finite state set $P$, a deterministic
binary machine $R$ with $K+(k+2)$ work tapes, and an exit state in $P$ such that the
following holds. For each $\alpha$, there are natural numbers $T$ and $B$ so that, for
every binary input $w$, after exactly $T$ steps on $\langle\alpha,w\rangle$, $R$ is in
the exit state, its input-head position is $1$, all work-head positions are $0$, and its
output is empty; work tape $k$ stores the decoded machine's transition table, work tape
$k+1$ stores its unary initial-state encoding shifted by two cells, and every other
work-tape cell is blank. The exit state is not reached at any earlier time, the output
remains empty through time $T$, and the total space used by time $T$ is at most $B$.
\end{theorem}

\begin{theorem}[Self-simulation with time and space bounds]
\lean{Turing.exists_selfUniversalK}
\label{Turing.exists_selfUniversalK}
\leanok
\uses{Turing.EffectiveCodeK, Turing.EffectiveCodeK.toFinTM, Turing.FinTM, Turing.FinTM.ComputesInTime, Turing.MultiTapeTM.initCfg, Turing.MultiTapeTM.spaceUsed, Turing.pairEncode}
For every natural number $k$ and every effective coding system that assigns a
$k$-work-tape Boolean machine with a finite state set to each Boolean string $\alpha$,
there is a finite-state deterministic Boolean machine $U$ such that, for each $\alpha$,
some positive integer $C$ works for every Boolean string $w$. Whenever the machine
assigned to $\alpha$ produces output $o$ on $\langle\alpha,w\rangle$ and halts at
exactly time $t$, $U$ produces the same output on that input at exactly time $C(t+1)$,
using at most $C(s+1)$ work-tape cells if the assigned machine uses exactly $s$ cells by
time $t$. Conversely, every output that $U$ produces on $\langle\alpha,w\rangle$ is
produced by the machine assigned to $\alpha$ after some finite number of steps.
\end{theorem}

\begin{theorem}[Universal simulation for a fixed number of tapes]
\lean{Turing.universalK}
\label{Turing.universalK}
\leanok
\uses{Turing.EffectiveCodeK, Turing.EffectiveCodeK.toFinTM, Turing.FinTM, Turing.FinTM.ComputesInTime, Turing.pairEncode}
For every natural number $k$ and every effective coding scheme $E$ for machines with $k$
work tapes, there is a finite binary-alphabet Turing machine $U$ such that, for each
code $\alpha$, there is a natural number $C$ for which the following holds for every
Boolean input word $x$. If the machine decoded from $\alpha$ produces output $y$ and
halts on $x$ after exactly $t$ steps, then $U$ produces $y$ and halts on
$\langle\alpha,x\rangle$ after exactly $C(t+1)$ steps; conversely, if $U$ produces $y$
and halts on $\langle\alpha,x\rangle$, then the machine decoded from $\alpha$ produces
$y$ and halts on $x$ after some number of steps.
\end{theorem}
